% =========================================================
% Compléments M1 — Lemme d'Itô (B. Théorie)
% =========================================================

\subsection*{(1) Forme intégrale : ce que la formule dit vraiment}
La version différentielle est très pratique, mais pour raisonner sans ambiguïté il est utile de retenir la \textbf{forme intégrale}.
Si $X$ suit
\[
X_t=X_0+\int_0^t \mu(s,X_s)\,\dd s+\int_0^t \sigma(s,X_s)\,\dd W_s
\]
et si $f\in C^{1,2}$, alors
\[
\boxed{
\begin{aligned}
f(t,X_t)=\;&f(0,X_0)
 +\int_0^t \partial_t f(s,X_s)\,\dd s
 +\int_0^t \partial_x f(s,X_s)\,\dd X_s\\
&+\frac12\int_0^t \partial_{xx} f(s,X_s)\,\dd [X]_s,
\end{aligned}}
\]
avec $[X]_s=\int_0^s \sigma(u,X_u)^2\,\dd u$.
En remplaçant $\dd X_s=\mu\,\dd s+\sigma\,\dd W_s$ et $\dd[X]_s=\sigma^2\,\dd s$, on obtient la forme classique :
\[
f(t,X_t)=f(0,X_0)
\!+\!\int_0^t\!\Big(\partial_t f+\mu\,\partial_x f+\tfrac12\sigma^2\partial_{xx} f\Big)(s,X_s)\,\dd s
\!+\!\int_0^t\!\sigma\,\partial_x f(s,X_s)\,\dd W_s.
\]
\paragraph{Lecture martingale.}
Le dernier terme est une intégrale d'Itô, donc une martingale (sous conditions). Le reste est un terme de drift (déterministe conditionnellement à la trajectoire).

\subsection*{(2) Pourquoi $C^{1,2}$ ? et que faire quand ce n'est pas vrai}
\paragraph{Rôle de $C^{1,2}$.}
Le développement de Taylor utilisé dans la preuve requiert une dérivée en temps et deux dérivées en espace pour contrôler les restes.

\paragraph{Payoffs non lisses en finance.}
Les payoffs vanilles ont des singularités :
\begin{itemize}[leftmargin=*]
\item Call : $\Phi(S)=(S-K)^+$ non $C^1$ en $S=K$.
\item Digital : $\Phi(S)=\1_{\{S>K\}}$ discontinu.
\end{itemize}
On contourne en pratique via :
\begin{itemize}[leftmargin=*]
\item PDE en sens faible / viscosité ;
\item approximation lisse (mollification) + passage à la limite ;
\item extensions (Tanaka, temps local) pour traiter des convexités concentrées.
\end{itemize}

\subsection*{(3) Itô multi-dimensionnel : notation propre et cas corrélé}
Soit $\bm{X}_t\in\R^d$ avec
\[
\dd \bm{X}_t=\bm{\mu}(t,\bm{X}_t)\,\dd t+\Sigma(t,\bm{X}_t)\,\dd \bm{W}_t,
\]
où $\bm{W}$ est un brownien $d$-dimensionnel à composantes indépendantes.
Alors, pour $f\in C^{1,2}$,
\[
\dd f(t,\bm{X}_t)
=\partial_t f\,\dd t+\nabla f^\top \dd \bm{X}_t+\frac12\mathrm{Tr}\!\big(\Sigma\Sigma^\top \nabla^2 f\big)\dd t.
\]
\paragraph{Cas corrélé.}
Si l'on part d'un brownien corrélé $\bm{\widetilde W}$ tel que $\dd \bm{\widetilde W}=L\,\dd \bm{W}$ avec $LL^\top=\rho$, on peut absorber la corrélation dans $\Sigma$.
Dans les modèles multi-facteurs, l'objet important est toujours $\Sigma\Sigma^\top$ : c'est lui qui porte les covariations.

\subsection*{(4) Itô--Tanaka (optionnel mais très utile pour les payoffs à kink)}
Pour une fonction convexe mais non $C^2$ (comme $x\mapsto |x|$), la formule d'Itô se prolonge via la \textbf{formule de Tanaka}.
Pour un semimartingale continu $X$ :
\[
|X_t|
=|X_0|+\int_0^t \mathrm{sgn}(X_s)\,\dd X_s + L_t^0(X),
\]
où $L_t^0(X)$ est le \emph{temps local} en 0.
\paragraph{Interprétation.}
Le temps local mesure combien de temps (au sens «densité de passage») la trajectoire passe près d'un niveau.
En finance, il apparaît naturellement quand on manipule des payoffs avec un kink (call/put) ou des barrières (densité de premier passage).
\paragraph{Message pédagogique.}
Même si l'on ne l'utilise pas directement pour pricer une vanilla, Tanaka explique pourquoi la convexité concentrée en $K$ fait apparaître des effets de couverture (Gamma) «localisés».

\subsection*{(5) Itô pour semimartingales : la version «générique» (message M1)}
Dans un cours de calcul stochastique, on retient que l'objet central est la semimartingale continue
\[
X_t=X_0+A_t+M_t,
\]
où $A$ est de variation finie et $M$ une martingale locale continue.
Alors, pour $f\in C^2$ (sans dépendance en $t$ pour simplifier),
\[
f(X_t)=f(X_0)+\int_0^t f'(X_s)\,\dd X_s+\frac12\int_0^t f''(X_s)\,\dd [X]_s.
\]
Cette forme est souvent la plus robuste conceptuellement : elle fait apparaître explicitement $\dd[X]_s$, donc l'endroit exact où la rugosité (variation quadratique) intervient.
